\section{Introduction}

The convergence of edge, fog, cloud, and high-performance computing (HPC) infrastructures has fundamentally transformed the design of eScience workflows. This integrated ecosystem, often termed the \emph{Compute Continuum}, enables researchers to distribute complex analytical pipelines across a hierarchy of resources. By placing latency-sensitive acquisition tasks near physical instruments and offloading heavy computational stages to the cloud or HPC centers, workflows can theoretically achieve both responsiveness and scale. However, this architectural shift introduces a significant scheduling bottleneck: the transition from centralized execution to a distributed, cross-layer environment turns every task placement into a high-stakes trade-off between compute capacity, data locality, and network availability.

Scheduling data-intensive directed acyclic graph (DAG) workflows across the continuum is particularly challenging because of the inherent coupling between task execution and data movement. In traditional cloud environments, high-bandwidth internal networks often allow schedulers to prioritize compute-load balancing. In the Compute Continuum, however, intermediate data must frequently traverse heterogeneous links with highly variable bandwidths. A placement that appears computationally optimal can easily violate strict scientific deadlines if the communication delay across a constrained link exceeds the time saved during execution. This problem is compounded in infrastructure-constrained settings, such as disaster management or real-time environmental monitoring, where the very events requiring rapid analysis often coincide with network instability or resource volatility.

Furthermore, the Compute Continuum is characterized by deep heterogeneity and unpredictable resource availability. Federated fog tiers, which bridge the gap between local edge devices and distant cloud centers, offer necessary elasticity but introduce complex inter-dependencies. Task execution times vary across diverse hardware, and the availability of these resources can fluctuate due to localized failures or competing workloads. Existing scheduling models often struggle to account for these dynamics simultaneously. Specifically, there is a lack of frameworks that can jointly optimize task placement and execution fractions while remaining resilient to the communication bottlenecks and node failures that define real-world distributed environments.

We address these challenges by proposing a data- and bandwidth-aware optimization framework designed for resilient DAG execution across the Compute Continuum. Our approach explicitly models the interplay between task-level data sizes and link-level bandwidth, enabling an orchestrator to evaluate placements based on the true combined cost of computation and communication. By integrating a mixed-integer linear programming (MILP) benchmark with scalable scheduling algorithms, we provide a foundation for managing time-critical eScience workflows that must remain performant even under resource failures and network constraints. This work focuses on establishing the theoretical and experimental evidence needed to ensure scientific workflows can navigate the complexities of modern distributed infrastructures without compromising on deadlines or reproducibility.

\subsection{Compute-Continuum Workflow Motivation}

The execution model considered in this study distinguishes three complementary roles. (1) Edge devices and sensors generate streams, and perform immediate acquisition. (2) Federated fog nodes run deadline-sensitive DAG tasks, coordinate nearby resources, and transfer intermediate data only when the resulting latency is justified. (3) Cloud and HPC systems provide elastic capacity for global optimisation, large simulation campaigns, archival processing, and training workloads that are not suitable for the constrained fog tier. This separation does not require every workflow task to traverse every tier; instead, it allows the scheduler to make placement decisions based on the data location, task deadline, resource capability, and inter-tier communication cost.

The design is relevant to high-performance scientific workflows because scientific pipelines frequently combine short control-loop stages with larger batch or ensemble stages. The same application may therefore need to preserve low-latency response locally while sending selected checkpoints, derived features, or nonurgent tasks to more powerful resources. The proposed DAG model captures these dependencies explicitly. Data size on each DAG edge and bandwidth on every resource link determine the communication component of a candidate placement. This view makes the federated fog tier an interoperable component of a broader Compute Continuum rather than a standalone deployment.

\begin{figure*}[t]
\centering
\begin{tikzpicture}[
  scale=0.94, transform shape,
  font=\small,
  box/.style={draw, rounded corners=2pt, align=center, minimum height=8mm, inner sep=4pt},
  tier/.style={box, fill=black!4, minimum width=29mm, minimum height=15mm},
  ctrl/.style={box, fill=black!8, minimum width=42mm},
  store/.style={box, fill=black!2, minimum width=41mm, minimum height=12mm},
  ar/.style={-{Latex[length=2mm]}, thick},
  dbl/.style={{Latex[length=2mm]}-{Latex[length=2mm]}, thick, dashed}
]
\node[ctrl] (orc) {\textbf{Workflow Orchestrator}};
\node[ctrl, below=7mm of orc] (sch)
  {\textbf{Scheduler}\\[1pt]\footnotesize MILP $\vert$ LP relaxation $\vert$ greedy $\vert$ GA};
\draw[ar] (orc) -- node[right=1mm, font=\scriptsize, align=left]
  {workflow descriptor:\\ $T$, edges $(T_j,T_i)$, $D_{ji}$, $t_{\max}$} (sch);

\node[tier, below left=13mm and 16mm of sch] (edge)
  {\textbf{Edge tier}\\\footnotesize data capture,\\\footnotesize initial processing};
\node[tier, below=13mm of sch] (fog)
  {\textbf{Federated Fog tier}\\\footnotesize \emph{modelled and}\\\footnotesize \emph{evaluated here}};
\node[tier, below right=13mm and 16mm of sch] (cloud)
  {\textbf{Cloud / HPC tier}\\\footnotesize large-scale storage,\\\footnotesize HPC analytics};

\draw[ar] (sch) -- node[left=1mm, font=\scriptsize] {$e_{ia}, l_{ia}$} (fog);
\draw[ar] (sch.south west) -- (edge.north);
\draw[ar] (sch.south east) -- (cloud.north);
\draw[dbl] (edge) -- node[below, font=\scriptsize] {$c_{ji}^{ba}=D_{ji}/B_{ba}$} (fog);
\draw[dbl] (fog)  -- node[below, font=\scriptsize] {$c_{ji}^{ba}=D_{ji}/B_{ba}$} (cloud);
\begin{scope}[on background layer]
  \node[draw, dashed, thick, rounded corners=3pt, fit=(fog), inner sep=3pt] {};
\end{scope}

\node[store, below =14mm of edge] (meta)
  {\textbf{Metadata Repository}\\[1pt]\footnotesize workflow + resource descriptors;\\
   \footnotesize $D_{ji}$, $B_{ba}$, $t_{ia}$, $A_i$, $P_{ia}$};
\node[store, below=14mm of fog] (prov)
  {\textbf{Provenance Store}\\[1pt]\footnotesize run manifests: seed, placements,\\
   \footnotesize metrics, hashes (RO-Crate)};
\draw[ar] (meta.north) -- node[left=1mm, font=\scriptsize] {descriptors} (edge.south);
\draw[ar] (fog.south) -- node[right=1mm, font=\scriptsize] {run record} (prov.north);
\draw[dbl] (meta) -- (prov);
\draw[ar] (meta.west) -- ++(-11mm,0) coordinate (fb) |- (orc.west);
\node[left=0.5mm of fb, font=\scriptsize] {state};
\end{tikzpicture}
\caption{Positioning of the proposed scheduling model within a Compute-Continuum workflow
architecture. The Orchestrator resolves a portable workflow descriptor together with current
resource state; the Scheduler component is where the MILP, randomized relaxation, greedy and
genetic methods of Sections~\ref{sec:system_model} and~\ref{sec:algorithms} are instantiated.
Every tier crossing is priced by the same term $c_{ji}^{ba}=D_{ji}/B_{ba}$, so a remote tier is
selected only when its capability advantage exceeds its transfer cost. The Metadata Repository
and Provenance Store realise the descriptor separation and run-level provenance of
Section~\ref{sec:reproducibility}. \emph{The dashed region indicates the tier modelled and
evaluated in Section~\ref{performance_analysis}}; cloud and HPC tiers are represented in the
model through the values of $t_{ia}$ and $B_{ba}$ rather than as distinct structural constructs.}
\label{fig:compute_continuum_architecture}
\end{figure*}

\subsection{Workflow Model and Research Questions}

The motivating application is represented as a DAG in which a node is a task and an edge denotes both a precedence relationship and the data transferred between tasks. The central research question is how to assign tasks to heterogeneous federated resources such that the workflow respects deadlines and communication constraints while maximising achieved accuracy. The study therefore addresses four linked questions: how data size and bandwidth change task placement; how a global optimisation model compares with scalable scheduling alternatives; how workflow scale and resource scale affect accuracy; and how the methods respond to workflow-node failures.

Figure~\ref{fig:compute_continuum_architecture} shows the latency-sensitive portion of the execution environment. The scheduler receives task and resource metadata, including task execution requirements, data-dependency sizes, link bandwidths, and deadline requirements. This metadata-first representation also supports the reproducibility and interoperability practices described later in the paper. It creates a clear boundary between a portable workflow description and the dynamic resource state used for a specific execution.


\subsection{Contributions}

This paper makes the following contributions:
\begin{itemize}[leftmargin=*]
    \item It positions data- and bandwidth-aware DAG scheduling as a Compute-Continuum workflow problem in which federated fog provides the latency-sensitive execution tier and cloud/HPC resources provide higher-capacity complementary tiers.
    \item It formulates task placement and approximate execution as a mixed integer linear program (MILP) that jointly captures task precedence, heterogeneous execution time, data-transfer size, inter-node bandwidth, permitted placements, and a workflow deadline.
    \item It retains and evaluates a randomized relaxation, a band\-width/data-aware greedy scheduler, and a genetic-algorithm scheduler as scalable alternatives to exact optimisation.
    \item It provides a four-dimensional performance comparison covering deadline flexibility, federated resource scale, DAG scale, and workflow-node failure resilience.
    \item It specifies a FAIR-oriented reproducibility methodology based on portable workflow metadata, versioned configurations, input catalogs, output records, and provenance capture, without claiming release of artifacts that are not yet publicly available.
\end{itemize}

\subsection{Relationship to the Conference Version and Novelty of This Submission}
\label{subsec:novelty}

The combinatorial optimisation machinery in this paper is prior work. The MILP
formulation, its correctness proofs and linearisation, the randomized relaxation, and the greedy
and genetic schedulers were introduced in our IEEE Cloud Summit paper~\cite{gupta2025efficient}
and are reproduced here unchanged. What is new in this submission is the level at which the model
is posed: task placement is reframed as an orchestration decision over a federated, multi-tier
continuum, and the model's parameters are shown to constitute a portable, FAIR-compatible workflow
description rather than solver input alone. Table~\ref{tab:novelty} states this division explicitly
so that the boundary between inherited and new material is not left to inference.

\begin{table}[t]
\centering
\caption{Delineation of material inherited from the conference version and material introduced in
this submission.}
\label{tab:novelty}
\small
\begin{tabular}{p{0.44\linewidth}p{0.48\linewidth}}
\toprule
\textbf{Inherited from the conference version} & \textbf{New in this submission} \\
\midrule
MILP formulation, correctness proofs (Lemmas~1--6, feasibility theorem), and linearisation &
Compute-Continuum framing: placement posed as orchestration across federated edge, fog, cloud and
HPC tiers rather than isolated fog scheduling (Section~\ref{sec:orchestration}) \\
\addlinespace
LP relaxation with randomized rounding &
Scheduler-selection policy binding each method to an operational condition and a provenance record
(Table~\ref{tab:orchestration_policy}) \\
\addlinespace
Bandwidth- and data-aware greedy heuristic &
Adaptive execution and resilience loop preserving the non-preemptive assumption under resource
failure and bandwidth degradation \\
\addlinespace
Genetic-algorithm metaheuristic &
Streaming and in-transit interpretation: per-window re-estimation of $D_{ji}$ and $B_{ba}$, with
transformation tasks made explicit on the data path \\
\addlinespace
Communication model $c_{ji}^{ba}=D_{ji}/B_{ba}$ and the accuracy--deadline trade-off via execution
fraction $l_{ia}$ &
Separation of a portable workflow descriptor from a deployment-specific resource descriptor, and
the FAIR-principle mapping of $D_{ji}$, $B_{ba}$, $t_{ia}$, $A_i$, $P_{ia}$
(Section~\ref{sec:reproducibility}, Table~\ref{tab:fair_mapping}) \\
\addlinespace
Simulation results across slack factor, resource scale, DAG scale and node failure &
Related-work positioning against Compute-Continuum orchestrators, scientific workflow engines and
reproducibility frameworks (Section~\ref{sec:related_work}); energy, cost and security expressed as
resource-descriptor metadata rather than informal assumptions \\
\bottomrule
\end{tabular}
\end{table}

The contribution addresses several themes of this Special Issue directly: orchestration and
resource-aware scheduling (Section~\ref{sec:orchestration}), resilience and adaptive execution
(Sections~\ref{sec:orchestration} and~\ref{performance_analysis}), data streaming and in-transit
processing (Section~\ref{sec:orchestration}), and workflow description and reproducibility
(Section~\ref{sec:reproducibility}).

The remainder of the paper reviews related work, defines the system model and optimisation formulation, describes the reproducibility methodology, presents the retained scheduling algorithms, evaluates the reported simulation results, and provides the original formulation proofs in an appendix.
